\documentclass{article}

\usepackage{mla13}

% ========== Note to users =========
% Remove these two packages after 
% you remove the example text.
\usepackage{blindtext}
\usepackage{lipsum}

\title{Example Paper Title}	% Title for the paper
\firstname{John}			% Author's first name
\lastname{Smith}			% Author's last name
\professor{Dr. Doctore}		% Last Name of the professor
\class{Example Class}		% Name of the class
%\date{10 April 2017}		% Uncomment to change due-date
\sources{references.bib}	% File holding all the sources

\begin{document}
% First create the title page on its own.
\maketitlepage
\newpage

% Then create the table of contents on its own.
\tableofcontents
\newpage

% Finally create the abstract on its own.
\begin{abstract}
	\blindtext[3] % Replace this line with the abstract.
\end{abstract}
\newpage

% Paper content begins here.
\makeheader	% Adds the MLA-Style information
\maketitle	% Redefined to use MLA-Style

% Put paper here.
% Begin example text.
This is here just so I can cite this thingy \cite{Kavanaugh1976} and also this thingy \cite{Saussure1995} so there can be some examples of bibliography and such. Usually this is where you'd put your introduction or something.
\blinddocument
% End example text.

% Put the works-cited on a separate page.
\newpage
\makeworkscited

\end{document}